% =====================================================================
% DISRetrieval Readiness Program � Status Report
% Compile: pdflatex report.tex  (twice for references)
% =====================================================================
\documentclass[11pt,a4paper]{article}

\usepackage[T1]{fontenc}
\usepackage[utf8]{inputenc}
\usepackage{lmodern}
\usepackage[a4paper,margin=2.3cm]{geometry}
\usepackage{parskip}
\usepackage{microtype}

\usepackage{amsmath,amssymb}
\usepackage{booktabs}
\usepackage{longtable}
\usepackage{tabularx}
\usepackage{array}
\usepackage{enumitem}
\setlist[itemize]{leftmargin=1.4em,itemsep=2pt,topsep=2pt}
\setlist[enumerate]{leftmargin=1.6em,itemsep=2pt,topsep=2pt}

\usepackage{xcolor}
\usepackage[most]{tcolorbox}
\newtcolorbox{decisionbox}{
  colback=blue!4, colframe=blue!55!black,
  title={\bfseries Decision}, fonttitle=\bfseries,
  boxrule=0.6pt, arc=2pt, left=6pt, right=6pt, top=4pt, bottom=4pt
}
\newtcolorbox{blockerbox}{
  colback=red!4, colframe=red!60!black,
  title={\bfseries Blocker}, fonttitle=\bfseries,
  boxrule=0.6pt, arc=2pt, left=6pt, right=6pt, top=4pt, bottom=4pt
}
\newtcolorbox{notebox}{
  colback=gray!5, colframe=gray!55!black,
  boxrule=0.4pt, arc=2pt, left=6pt, right=6pt, top=4pt, bottom=4pt
}

\usepackage{fancyhdr}
\pagestyle{fancy}
\fancyhf{}
\fancyhead[L]{\small RAG Architecture Portfolio --- Status Report}
\fancyhead[R]{\small\thepage}
\renewcommand{\headrulewidth}{0.3pt}

\usepackage[hidelinks]{hyperref}
\hypersetup{
  pdftitle={RAG Architecture Portfolio Status Report},
  pdfauthor={Project Owner},
  bookmarksnumbered=true
}

\usepackage{titlesec}
\titleformat{\section}{\Large\bfseries}{\thesection}{0.6em}{}
\titleformat{\subsection}{\large\bfseries}{\thesubsection}{0.5em}{}
\titlespacing*{\section}{0pt}{1.2em}{0.4em}
\titlespacing*{\subsection}{0pt}{0.9em}{0.3em}

\newcommand{\code}[1]{\texttt{#1}}

% =====================================================================
\begin{document}

\begin{center}
{\LARGE\bfseries RAG Architecture Portfolio}\\[4pt]
{\large Status Report and Composition Decision}\\[6pt]
\rule{\linewidth}{0.4pt}
\end{center}

\vspace{0.5em}

\begin{notebox}
\textbf{Scope.} This document records the architecture-level decisions
taken during the preparation of a controlled retrieval comparison
program. It states what was verified, what was excluded, and why, in
order to make the final comparison defensible under academic review.
No empirical retrieval result is reported here.
\end{notebox}

% =====================================================================
\section{Purpose and Objective}

The program aims to answer a single question:

\begin{quote}
\emph{Which retrieval architecture recovers relevant information most
effectively from long-context corpora?}
\end{quote}

The question is retrieval-only. Final answer generation, LLM-based
judging, and answer evaluation are disabled across every condition.
The unit of comparison is the architecture-defining retrieval logic.

The program operates under three binding constraints that shape every
decision below:

\begin{enumerate}
  \item \textbf{Academic correctness} --- every claim must be
    evidence-based, reproducible, and defensible.
  \item \textbf{Comparative fairness} --- all included architectures
    must be evaluated under common semantic representation, common
    datasets, common queries, and a common evaluation contract.
  \item \textbf{Architectural integrity} --- no architecture-defining
    component may be silently substituted. Where substitution is
    unavoidable, it is labelled and its risk is stated.
\end{enumerate}

The program also operates under two hard resource constraints:

\begin{itemize}
  \item \textbf{Local execution only.} No hosted LLM API is available.
    Any architecture that requires a hosted LLM at index time or query
    time is out of scope.
  \item \textbf{CPU-bound staging.} GPU compute is limited to a single
    Colab T4 session and is reserved for embedding generation only.
\end{itemize}

% =====================================================================
\section{DISRetrieval: Terminal Disposition}

\begin{blockerbox}
DISRetrieval is closed as \code{DEFERRED\_WITH\_REASON}.\\
Reason: \code{BLOCKED\_ON\_ARCHITECTURE\_DEFINING\_RST\_PARSER\_LICENSE}.
\end{blockerbox}

\subsection{Verification findings}

The vendor repository at commit
\code{4f3144042c6eb1e514a0a89ec704d892a66e700c} was audited
statically. The following facts were established:

\begin{itemize}
  \item \textbf{No license file is present.} No \code{LICENSE},
    \code{NOTICE}, or \code{COPYING} file exists. No per-file license
    header exists. Under default copyright, execution, reuse,
    modification, and redistribution are all prohibited.
  \item \textbf{The RST parser is architecture-defining.} The
    discourse-aware hierarchical representation produced by the
    neural RST parser is the paper's primary contribution. Substituting
    the parser with a heuristic, a generic constituency parser, or a
    clean-room reimplementation would materially alter the
    architecture.
  \item \textbf{The neural parser ships with a bundled checkpoint.}
    The trained model weights are included in the repository. Loading
    them constitutes a use of unlicensed weights.
  \item \textbf{The bundled checkpoint is English-only.} No
    Chinese-compatible checkpoint is provided. The bundled model is
    trained on RST-DT and is not suitable for Chinese datasets.
\end{itemize}

\subsection{Binding decisions}

\begin{decisionbox}
\begin{tabularx}{\linewidth}{@{}l X@{}}
\toprule
Q1 (RST parser policy)     & \textbf{A} --- defer entire architecture. No substitution. \\
Q2 (track placement)       & \textbf{DEFER} --- remains in Track C appendix. \\
Q3 (\code{multiFieldQA-zh}) & \textbf{EXCLUDE} --- no language-compatible checkpoint. \\
Q4 (institutional license) & \textbf{NOT ESTABLISHED} --- no research exception evidence. \\
\bottomrule
\end{tabularx}
\end{decisionbox}

\subsection{Reopening conditions}

The architecture may be reconsidered if and only if one of the
following occurs:

\begin{itemize}
  \item The vendor repository adds a compatible open-source license.
  \item The repository owner issues written permission for execution.
  \item The institution issues written, specific legal clearance for
    the exact commit, the exact checkpoint SHA-256, and the permitted
    research scope.
  \item A legally licensed equivalent parser and checkpoint are
    identified.
\end{itemize}

Reopening requires a new versioned amendment and a fresh synthetic
acceptance run before any dataset pilot.

% =====================================================================
\section{Excluded Architectures}

The following architectures were excluded from the primary comparison.
Exclusion is a methodological choice, not a failure of the comparison.
It is recorded here to make the final scope defensible.

\subsection*{RAPTOR}

\begin{itemize}
  \item \textbf{Blocker:} index-time LLM summarisation is
    architecture-defining and cannot be removed.
  \item \textbf{Cost:} a single dataset requires thousands of LLM
    summarisation calls. The cost is not sustainable without hosted
    infrastructure.
  \item \textbf{Decision:} excluded from the primary comparison.
\end{itemize}

\subsection*{HippoRAG2}

\begin{itemize}
  \item \textbf{Blocker:} index-time OpenIE (named-entity and triple
    extraction) is architecture-defining.
  \item \textbf{Cost:} graph construction requires on the order of
    \(10^{7}\) LLM tokens for a single long-context dataset.
  \item \textbf{Decision:} excluded from the primary comparison.
\end{itemize}

\subsection*{LightRAG}

\begin{itemize}
  \item \textbf{Blocker:} index-time entity and relation extraction is
    architecture-defining.
  \item \textbf{Cost:} graph construction is the most expensive of the
    three graph-based systems considered, exceeding \(10^{7}\) LLM
    tokens per dataset.
  \item \textbf{Decision:} excluded from the primary comparison.
\end{itemize}

\subsection*{MemorAG}

\begin{itemize}
  \item \textbf{Blocker:} requires a GPU with 16--24~GB VRAM plus an
    LLM at both index time and query time.
  \item \textbf{Decision:} excluded. Hardware requirements exceed the
    available host envelope.
\end{itemize}

\subsection*{KAG}

\begin{itemize}
  \item \textbf{Blocker:} requires a multi-service environment
    (OpenSPG, Neo4j, hosted LLM).
  \item \textbf{Decision:} excluded. The infrastructure is not
    reproducible within the scope of the study.
\end{itemize}

\subsection*{HiKEY}

\begin{itemize}
  \item \textbf{Blocker:} multimodal input contract required
    (PDF, tables, images, layout). No public repository is available.
  \item \textbf{Decision:} excluded. The canonical envelope is
    text-only.
\end{itemize}

\subsection*{Summary}

\begin{center}
\begin{tabularx}{\linewidth}{@{}l X l@{}}
\toprule
\textbf{Architecture} & \textbf{Primary blocker} & \textbf{Track} \\
\midrule
DISRetrieval & License absent on architecture-defining RST parser & C \\
RAPTOR       & Index-time LLM summarisation                   & C \\
HippoRAG2    & Index-time LLM OpenIE                          & C \\
LightRAG     & Index-time LLM entity/relation extraction      & C \\
MemorAG      & GPU + LLM requirements                         & C \\
KAG          & Multi-service environment                      & C \\
HiKEY        & Multimodal input contract                      & C \\
\bottomrule
\end{tabularx}
\end{center}

% =====================================================================
\section{On the Correctness of Scope Reduction}

Reducing the architecture set is sometimes perceived as a weakness. It
is not. Under the constraints stated in Section~1, the following
outcomes hold.

\begin{itemize}
  \item \textbf{Including a system that cannot be run fairly is worse
    than excluding it.} A partial or substituted execution would
    violate the fairness pillar and produce numbers that cannot be
    defended.
  \item \textbf{Exclusion with a stated reason is reproducible.} Any
    reviewer can re-derive the exclusion from the same repository
    audit, the same cost estimates, and the same license findings.
  \item \textbf{Exclusion increases internal validity.} The remaining
    architectures share the same canonical envelope, the same
    embedder, the same queries, the same metrics, and the same
    evaluation contract. The comparison is therefore clean.
\end{itemize}

A feasibility appendix that records the excluded systems is a
recognised methodological practice. It does not weaken the primary
result; it bounds the claims.

% =====================================================================
\section{Final Composition}

The primary comparison will include four architectures. All four are
runnable locally without hosted LLM calls and without GPU inference
beyond a single Colab T4 session used only for embedding generation.

\begin{center}
\begin{tabularx}{\linewidth}{@{}c l X l@{}}
\toprule
\textbf{\#} & \textbf{Architecture} & \textbf{Retrieval paradigm} & \textbf{State} \\
\midrule
1 & \code{my\_rag}     & Hybrid hierarchical (dense + sparse, parent aggregation) & Ready \\
2 & \code{longrag}     & Dense long-unit retrieval                                & Ready \\
3 & \code{pathrag}     & Path-based resource propagation                          & Partial \\
4 & \code{metadata\_rag} & Metadata-filtered hybrid retrieval                     & Planned \\
\bottomrule
\end{tabularx}
\end{center}

\subsection*{Notes on individual entries}

\paragraph{\code{my\_rag}.} Project-owned reference. Fully implemented
and validated in the infrastructure pilot. Serves as the anchor
architecture for the comparison.

\paragraph{\code{longrag}.} Controlled adaptation. Synthetic acceptance
passed. Storage equivalence validated against an exact reference
dense index. Embeddings are generated in Colab T4. No LLM stage is
required.

\paragraph{\code{pathrag}.} Project-owned paper-derived comparator.
Because the foreign repository is unlicensed, only the clean-room
comparator is executable. The comparator operates without embedding
augmentation. This is a structural limitation, not an algorithmic
inferiority, and it will be stated as such in every reported table.
Lower retrieval metrics are expected and are not treated as evidence
about the paper's original design.

\paragraph{\code{metadata\_rag}.} Paper-derived clean-room
implementation. No foreign source repository exists. Implementation
proceeds from the paper specification only, with a rule-based query
metadata inference stage that does not require an LLM.

% =====================================================================
\section{Datasets}

Four datasets will be used in the primary comparison. All four are
registered in the canonical suite and have been verified for
integrity.

\begin{center}
\begin{tabularx}{\linewidth}{@{}l r r X@{}}
\toprule
\textbf{Dataset} & \textbf{Documents} & \textbf{Queries} & \textbf{Primary evaluation unit} \\
\midrule
LaRA            & 131   & 2{,}326 & Canonical document \\
CLERK           & 2{,}756 & 1{,}000 & Passage \\
CUAD            & 510   & 20{,}910 & Source span \\
GreekBarBench   & 284   & 284     & Source span \\
\bottomrule
\end{tabularx}
\end{center}

The four datasets provide coverage along three dimensions:

\begin{itemize}
  \item \textbf{Domain:} legal, financial, general, and multilingual.
  \item \textbf{Language:} English and Greek.
  \item \textbf{Corpus scope:} global, shared-domain, and
    query-specific.
\end{itemize}

% =====================================================================
\section{What Is Not Claimed}

The following statements are explicitly out of scope for this report
and for the resulting primary comparison.

\begin{itemize}
  \item No claim of paper-level effectiveness is made for any excluded
    architecture.
  \item No numeric comparison is made between included architectures
    and excluded architectures.
  \item No comparison is made between the included architectures and
    any hosted or commercial retrieval system.
  \item No claim is made that the included set is a complete or
    representative sample of the retrieval literature.
  \item No claim is made that lower metrics observed in the
    \code{pathrag} comparator reflect the design of the original
    PathRAG system.
\end{itemize}

% =====================================================================
\section{Summary of Decisions}

\begin{center}
\begin{tabularx}{\linewidth}{@{}l X@{}}
\toprule
\textbf{Decision area} & \textbf{Resolution} \\
\midrule
DISRetrieval & Deferred with reason: architecture-defining RST parser license absent. \\
RAPTOR       & Excluded: index-time LLM cost not sustainable. \\
HippoRAG2    & Excluded: index-time LLM cost not sustainable. \\
LightRAG     & Excluded: index-time LLM cost not sustainable. \\
MemorAG      & Excluded: GPU and LLM requirements. \\
KAG          & Excluded: service environment required. \\
HiKEY        & Excluded: multimodal contract required. \\
Primary set  & Four architectures: \code{my\_rag}, \code{longrag}, \code{pathrag}, \code{metadata\_rag}. \\
Datasets     & LaRA, CLERK, CUAD, GreekBarBench. \\
Evaluation   & Retrieval-only, document or span level, common metrics. \\
LLM usage    & Zero hosted LLM calls in the primary comparison. \\
\bottomrule
\end{tabularx}
\end{center}

\vspace{1em}
\begin{center}
\rule{0.5\linewidth}{0.4pt}
\end{center}

\end{document}